\documentclass[11pt,a4paper]{article}
\usepackage[T1]{fontenc}
\usepackage{lmodern,microtype,amsmath,amssymb,amsthm,mathtools}
\usepackage[margin=29mm]{geometry}
\usepackage{booktabs,longtable,array,enumitem,xcolor}
\usepackage[hidelinks]{hyperref}
\usepackage{xurl}
\usepackage{fancyhdr}
\pagestyle{fancy}\fancyhf{}\fancyhead[L]{\small Theta-cusp causal scattering}\fancyhead[R]{\small M. Haufschild}\fancyfoot[C]{\thepage}
\setlength{\headheight}{14pt}
\newtheorem{theorem}{Theorem}[section]
\newtheorem{proposition}[theorem]{Proposition}
\newtheorem{lemma}[theorem]{Lemma}
\theoremstyle{definition}\newtheorem{definition}[theorem]{Definition}
\newtheorem{problem}[theorem]{Open problem}
\theoremstyle{remark}\newtheorem{remark}[theorem]{Remark}
\newcommand{\C}{\mathbb C}\newcommand{\R}{\mathbb R}\newcommand{\N}{\mathbb N}
\newcommand{\XiC}{\Xi_{\!c}}\newcommand{\RH}{\mathrm{RH}}
\newcommand{\FP}{\operatorname{FP}}\newcommand{\Res}{\operatorname{Res}}
\newcommand{\re}{\operatorname{Re}}\newcommand{\im}{\operatorname{Im}}
\newcommand{\dd}{\,\mathrm d}\newcommand{\id}{\mathrm{id}}
\newcommand{\cert}[1]{\par\smallskip\noindent\textit{Formal source:} \hyperref[cert:#1]{source group #1}.\par\smallskip}
\newcommand{\code}[1]{\ifmmode\text{\small\nolinkurl{#1}}\else{\small\nolinkurl{#1}}\fi}
\newcommand{\status}[1]{\textsc{#1}}
\allowdisplaybreaks
\setlist{nosep}
\title{Theta-cusp causal scattering and the Riemann hypothesis:\\formal results and unresolved interfaces of Route HB}
\author{Michael Haufschild}
\date{13 September 2026}
\begin{document}
\maketitle
\begin{abstract}
Route HB is a proposed connection between the Riemann hypothesis, automorphic cusp scattering, and nonlinear perturbations of a topological anti-de Sitter background. We give a source-audited account of its mathematical content and distinguish three different achievements: construction of an arithmetic resonance, nonvanishing of a specified forced response, and a dynamical law forcing a diagnostic to vanish. The first does not imply the second, and neither supplies the third. The formal corpus contains an entire spin-two source factorization, an exact five-factor criterion for survival of a fundamental forced pole, a universal onset-duration theorem, finite Laurent particular solutions, a nonseparated scalar--scalar sourced metric, and specified renormalized angular and radial identities. It also contains an explicit rational kernel whose real part detects displacement from the critical line above height 92. We present proofs of the principal identities and classify the missing identifications and hypotheses. Several stronger conclusions in working notes, including full spectral dominance and extinction of the complete physical action, are not certified results. No proof or disproof of the Riemann hypothesis is obtained. An accompanying declaration manifest and source snapshot distinguish formal statements from their intended physical applications.
\end{abstract}
\noindent\textbf{Keywords:} Riemann hypothesis; Eisenstein series; quasinormal modes; finite-part integrals; formalized mathematics; Lean.
\tableofcontents
\section{Scope, contribution, and evidence}

An arithmetic scattering pole is an occurrence statement. A proof of the Riemann hypothesis requires a restriction on where such poles may occur. Route HB attempts to obtain that restriction from a common time convention and the nonlinear action of a theta-derived source. Its proposed chain is
\[
\begin{gathered}
\text{theta source}\longrightarrow\text{Xi occurrence}
\longrightarrow\text{retained nonlinear source},\\
\text{retained nonlinear source}\dashrightarrow\text{complete physical causal action}
\dashrightarrow\text{extinction}\longrightarrow\RH.
\end{gathered}
\]
The dashed arrows are research objectives, not theorems. The final implication refers to the specified reduced extinction receivers, whose identification with a diagnostic of the complete action remains open. In particular, an equivalence between an extinction property and \(\RH\) does not establish either side.

\textbf{Viability status.} Route HB is not yet a validated proof strategy with only technical lemmas missing. Two substantive bridges remain unestablished: identification of the complete physical action with a faithful diagnostic, and an independent source law forcing the required diagnostic to vanish. The second is not a consequence of finishing the construction in the first. Either proposed bridge could fail. Any viable law must exclude hypothetical off-line zeros while remaining compatible with the known critical-line zeros; it need not, and cannot without further justification, demand that the entire physical action vanish. The formal components therefore do not yet establish the viability of the proposed RH mechanism.

This article is a critical exposition of a specified formal corpus, not a claim of a new zero-location theorem. Its contribution is to make the mathematical statements, their domains, and the unfinished transfers between them explicit. The results are organized around analytic factorization, response residues, sourced metric reconstruction, boundary renormalization, and the difference between a detector and a law that makes it vanish. The development is substantial, but the extent of a formal library is not evidence of progress on an unproved final implication.

The geometric setting is related to the exact topological-black-hole calculations of Aros, Mart\'inez, Troncoso, and Zanelli \cite{aros}, the gauge-invariant master equations of Kodama and Ishibashi \cite{ki}, and the gravitational quasinormal-frequency calculation of Birmingham and Mokhtari \cite{bm}. These works provide context for the choice of operators and boundary conditions. They do not identify the nonlinear arithmetic source with an extinction mechanism. The cusp integrals also require the distinction between ordinary integration and regularization familiar from Zagier's work \cite{zagier}.

\subsection{What is meant by a formal result}

A result labelled as a theorem below has an identified theorem declaration in the accompanying corpus manifest. The ordinary restrictions in its statement---for example, \(r>1\), a nonzero denominator, or membership in the type of actual Xi zeros---are part of its mathematical domain. No unproved analytic existence, positivity, extinction, or \(\RH\) premise is silently discharged. An implication with such a premise is explicitly called conditional, even if Lean formalizes the implication itself.

There are two distinct audits. A \emph{logical audit} reads the complete declaration type and the definitions and proof on which its mathematical interpretation depends. A \emph{kernel audit} executes the proof checker and records its transitive axioms. This manuscript's new work is a source and dependency audit. A build started during preparation was interrupted; no successful fresh kernel audit is claimed. The supplied \code{Verify.lean} is a reviewer verification driver, not an executed certificate from this preparation. Previous checks reported in research notes are provenance, not replacements for a fresh check of a changed source file. This distinction is retained in the supplementary review report.

Lean's usual foundational constants \code{propext}, \code{Classical.choice}, and \code{Quot.sound} are not conjectural mathematical hypotheses. Conversely, absence of an extra axiom does not remove a theorem's explicit arguments or assumptions hidden in a structure. Throughout, an abstract response satisfying prescribed properties is distinguished from the literal constructed response.

\subsection{Snapshot and limits of coverage}

The primary inputs were \code{docs/notes.md} and the Route HB entry of \code{rh-current-frontier-map.md}. The latter explicitly uses a 12 September 2026 evidence cutoff. Later source files close some tasks still described there as open; the notes also contain superseded and speculative passages. The accompanying inventory lists 1,250 HB-related Lean files and 22,138 declaration names. This is a lexical inventory, not a claim that every declaration is independently explained or checked here. Every item in the notes' numbered route ledger is addressed in the coverage table; the formal-source appendix identifies the specific statements used in the mathematical exposition.

Hybrid and deep corpus search were used for discovery, followed by exact source searches and reading. The search service reported broad source coverage but incomplete compiler metadata, with many candidates marked \code{source_only} and \code{applicability: not_checked}. Therefore search rank, a descriptive comment, and a ledger label such as ``FULL'' are never treated as proof certificates. Negative findings mean that a requisite theorem was not located and established in this audit, not that no differently named theorem can exist.

\section{Conventions and the two spectral parameterizations}

Write
\[
\Lambda(s)=\pi^{-s/2}\Gamma(s/2)\zeta(s),\qquad
\xi(s)=\tfrac12s(s-1)\Lambda(s),\qquad
\XiC(z)=\xi(\tfrac12+iz),
\]
where \(\xi\) is extended to an entire function. Equalities using \(\Lambda\) at its poles mean meromorphic identities, not ordinary finite values. The corpus defines \(\xi\) by an entire pole-cleared expression. A zero copy \(q\) consists of a zero \(\rho\) and one of its multiplicity labels. No assumption that zeros are simple is made. Set
\[
 z_q=-i(\rho-\tfrac12),\qquad
 \rho=\sigma+it,\qquad
 z_q=t+i(\tfrac12-\sigma).
\]
The critical-line condition is \(\im z_q=0\). Functional-equation reflection of a zero is \(\rho\mapsto1-\rho\); complex conjugation is a different operation.

\subsection{Optical scattering coordinate}

Define the normalized meromorphic quotient and its shadow by
\[
 \mathcal S(z)=\frac{\XiC(z+i)}{\XiC(z)},\quad
 j(z)=-z-i,\quad a(z)=\tfrac14+\tfrac i2z,\quad
 \nu(z)=\tfrac z2+\tfrac i4.
\]
Thus \(a=\tfrac12+i\nu\). The quotient is the corpus's normalized modular cusp coefficient; the classical unnormalized coefficient has an additional elementary factor. The frequencies used in the incidence packet are
\[
 \omega_\pm(n,z)=\pm\nu(z)-i(2n+\tfrac32).
\]
These formulas define the spectral data being audited. They do not assert an \(L^2\) harmonic at a continued arithmetic resonance, nor a nonlinear Einstein--Klein--Gordon realization of every such resonance.

\begin{proposition}[Optical incidence and shadow symmetry]\label{thm:incidence}
For every Xi zero copy \(q\) and \(n\in\N\), \(\mathcal S\) has negative meromorphic order at \(z_q\), is analytic and zero at \(j(z_q)\), and
\[
 \omega_+(n,j(z_q))=\omega_-(n,z_q).
\]
For every complex \(z\),
\[
 \im\omega_+(n,z)-\im\omega_-(n,z)=\tfrac12+\im z.
\]
Any real-valued reporter depending only on \(Q(z)=\tfrac14+\nu(z)^2\) has the same value at \(z\) and \(j(z)\).
\end{proposition}
\cert{A1}
\begin{proof}
All zeros of \(\XiC\) lie in \(|\im z|<1/2\). At \(z_q\), the numerator has argument of imaginary part in \((1/2,3/2)\), hence is nonzero. The denominator has positive zero order; their quotient has negative order. At the shadow, evenness of \(\XiC\) makes the numerator \(\XiC(-z_q)=0\), while its denominator is outside the zero strip. Moreover \(\nu(j(z))=-\nu(z)\). Substitution gives both the branch exchange and invariance of \(Q\). Finally the difference of imaginary parts is \(2\im\nu(z)=\im z+1/2\).
\end{proof}

The last assertion is a statement about reporters of the displayed form. It does not classify every optical observable. Requiring the splitting to equal \(1/2\) at all Xi copies is a zero-location requirement, not a consequence of its name ``Hawking splitting.''

\subsection{Metric-source rates and physical testing frequencies}

For the later metric-source construction define
\[
 s_+(\rho)=\rho-2,\qquad s_-(\rho)=-\rho-1,\qquad
 \lambda=\rho(\rho-1)=(s_\pm+1)(s_\pm+2).
\]
A quadratic source has time rate \(2s_\pm\), and its output angular shift is
\(\lambda_2=(2s_\pm+1)(2s_\pm+2)\).
The \emph{original optical testing frequencies} in the Laplace variable are instead
\begin{equation}\label{eq:testing}
 w_+=-1-\rho/2,\qquad w_-=-2+\rho/2.
\end{equation}
These are the formulas for \code{physicalFundamentalFrequency}. They are not \(s_\pm\). The same short name \code{physicalRate} is used in two namespaces for a parent rate and its doubled rate; the manifest uses full namespaces. The residue statements below use the parent rate of \code{D4BMModularCore}.
\cert{D1}
The distinction matters: a calculation at \(\mu=\lambda\) for an elementary metric mode cannot be inserted into an original optical left-reader pairing without an explicit matching theorem.

\subsection{Background and angular source}

Take exterior coordinates \((\tau,r,x,y)\), \(r>1\), \(y>0\), and the background
\[
 g_0=-(r^2-1)\dd\tau^2+\frac{\dd r^2}{r^2-1}
       +r^2\frac{\dd x^2+\dd y^2}{y^2}.
\]
This fixes the anti-de Sitter radius to one. Let \(\Delta=y^2(\partial_x^2+\partial_y^2)\); thus its \(L^2\) sign is nonpositive. The literal scalar Eisenstein field, in the corpus normalization, is denoted \(S_\rho\). It retains both Weyl terms and the entire Fourier--Bessel tail. Its constant terms are
\(\Lambda(2\rho-1)y^{1-\rho}+\Lambda(2\rho)y^\rho\), and its scalar equation is \(\Delta S_\rho=\lambda S_\rho\). The quadratic source uses \(S_\rho^2\), not \(|S_\rho|^2\). Analytic continuation and complex bilinear pairing are retained until a conjugation theorem is available. Local identities on \(y>0\), descent to a modular domain, and a renormalized integral over its cusp are different assertions.

\cert{D5}

\section{Canonical spin-two normalization}

The retained spin-two vertex has a useful entire factorization. This concerns a particular completed balanced vertex at the repeated Eisenstein parameter used in the corpus; it is not a nonvanishing theorem for every physical source functional.

\begin{theorem}[Entire amputation]\label{thm:amputation}
Let \(V(z)\) be the corpus's completed balanced spin-two vertex at repeated centered parameter \(3/2\), evaluated at projecting parameter \(iz\). Then
\begin{align}
 V(z)&=\XiC(z)^2U(z),\label{eq:amp}\\
 U(z)&=\frac{(4z^2+1)(4z^2+25)}{32}
          \XiC(z-3i)\XiC(-z-3i).\label{eq:unit}
\end{align}
The function \(U\) is entire and is the unique entire function satisfying \eqref{eq:amp}. In \(-5/2<\im z<5/2\), its zeros are precisely \(z=\pm i/2\). In particular it is nonzero on \(|\im z|<1/2\).
\end{theorem}
\cert{A2}
\begin{proof}
The completed-vertex factorization separates its projecting external square from the two shifted Xi factors and the displayed polynomial. Substitution of \(iz\) into that factorization gives \eqref{eq:amp}--\eqref{eq:unit}. Each factor in \eqref{eq:unit} is entire. In the larger open strip both shifted arguments have imaginary part less than \(-1/2\), so their Xi values are nonzero. The polynomial roots are \(\pm i/2\) and \(\pm5i/2\), with the latter two outside the open strip. For uniqueness, \(\XiC(0)\ne0\) gives a neighborhood on which cancellation of \(\XiC^2\) is valid. Any competing entire quotient agrees with \(U\) there and hence everywhere by analytic uniqueness.
\end{proof}

The external square has not disappeared from the source. Dividing by it is an analytic normalization supported by \eqref{eq:amp}; it is not permission to regard the unamputated vertex as nonzero at a Xi zero. Nor does uniqueness of this normalization construct a physical left testing field.

\section{The complete angular coefficient and forced-pole survival}

\subsection{Angular coefficient}

Denote by \(A_\rho\) the corpus's \code{physicalSelfSourceCoefficient}, obtained by assembling the residual angular period and all three principal cutoff corrections. Its normalization must be preserved: replacing that sum with only its residual part changes the following formula.

\begin{theorem}[Angular value and collision criterion]\label{thm:angular}
For a Xi zero copy with \(\re\rho>1/2\),
\begin{equation}\label{eq:fourlambda}
 A_\rho=\tfrac12\Lambda(5\rho/2-1)\Lambda(\rho/2)^2
                         \Lambda(1-3\rho/2).
\end{equation}
Moreover,
\begin{equation}\label{eq:collisions}
 A_\rho=0\quad\Longleftrightarrow\quad
 \xi(5\rho/2-1)=0\ \lor\ \xi(\rho/2)=0\ \lor\ \xi(3\rho/2)=0.
\end{equation}
\end{theorem}
\cert{A3}\cert{A4}
\begin{proof}
The three principal Mellin terms sum to the native principal moment. Each cutoff quotient has exactly the scalar factor needed to cancel the principal correction in the residual period. Substituting the evaluated residual period gives \eqref{eq:fourlambda}; this is the content of the source's cutoff-assembly identity, not an appeal to formal integration of a divergent series. All three affine arguments are nonreal because \(\im\rho\ne0\). Their Gamma and elementary pole-clearing factors are finite and nonzero, so vanishing of \(\Lambda\) is equivalent to vanishing of \(\xi\) at each argument. Applying \(\xi(1-u)=\xi(u)\) to the last factor gives \eqref{eq:collisions}.
\end{proof}

Equation \eqref{eq:fourlambda} evaluates a source; it does not prove that source nonzero. Its strict-right restriction is also essential to its present formal application. The value of this angular functional and the entire source \(U\) in Theorem~\ref{thm:amputation} are different objects.

\subsection{Onset and the finite-part radial moment}

Let \(\chi\) be the fixed smooth transition used in the corpus, zero before 0 and one after 1. Its derivative has integral one. Define
\[
 O(s,w)=\int_0^1\chi'(u)e^{(2s-w)u}\dd u,\qquad
 O_L(s,w)=\int_0^1\chi'(u)e^{L(2s-w)u}\dd u.
\]
These are \code{clearedOnset} and \code{clearedOnsetDuration}; the second formula stretches the exponent rather than silently replacing an already constructed solution by another one.

For \(0\le z<1\), define
\begin{align*}
 c_0(s)&=(5s^2+15s+4)/4,\\
 c_1(s,w)&=\frac{s(s-1)}4\left((s+2)^2+(s+1)(s+2)
                                  -\frac{(w+1)(w+2)}2\right),\\
 F(w,s,a;z)&=a^2\sqrt{1-z}\,(1+\sqrt{1-z})^{2s-w}
                      [c_0(s)+c_1(s,w)(1-z)].
\end{align*}
Powers of a positive real base use its real logarithm. Write \(b=1/2\). The radial moment used here is precisely
\begin{align}\label{eq:moment}
 M(w,s,a)={}&\frac{F(w,s,a;0)b^{w+1}}{w+1}
 +\int_0^b z^w[F(w,s,a;z)-F(w,s,a;0)]\dd z\nonumber\\
 &+\int_b^1z^wF(w,s,a;z)\dd z.
\end{align}
This is the source definition \code{fundamentalSourceMoment}. For the physical nonreal frequencies in \eqref{eq:testing}, \(\re w>-2\) and \(w\ne-1\); the endpoint subtraction gains the one power needed at zero. Formula \eqref{eq:moment} must not be replaced by an unsubtracted improper integral in that chamber.

\begin{theorem}[A common nonvanishing onset duration]\label{thm:onset}
There is one \(L>0\) such that \(O_L(s_b(\rho),w_e(\rho))\ne0\) for every Xi zero copy, both parent branches \(b\), and both testing sides \(e\). Also \(O(s,2s)=1\) and \(O_1(s,w)=O(s,w)\).
\end{theorem}
\cert{A6}\cert{A7}
\begin{proof}
The pulse transform \(z\mapsto\int_0^1\chi'(u)e^{zu}\dd u\) is entire and takes value one at zero. For each fixed pair \((s,w)\), its composition with \(L\mapsto L(2s-w)\) is entire and not identically zero. Its real nonzero set is open and dense; if \(2s=w\), it is the whole real line. Xi zero copies form a countable type, and the two branch choices are finite. Baire's theorem makes the intersection of these open dense sets dense. It meets \((0,\infty)\). The two identities follow by substitution and the normalization of \(\chi'\).
\end{proof}

This theorem establishes no preferred numerical duration and, in particular, does not certify \(L=1\) at every target frequency. Transporting the complete source, response, boundary terms, and time convention to the chosen \(L\) remains an additional obligation.

\subsection{Exact residue factorization}

Let \(R_{q,b}(w;a,x)\) denote \code{continuedPhysicalSelfResponse}, with tortoise observation point \(x<0\). Set \(z(x)=\tanh^2x\), and
\[
 m_w(x)=z(x)^{w/2}(1-z(x))^{1/2}.
\]
Let \(D'_{q,b,e}\) be \code{physicalFundamentalDenominatorSlope}. The corpus's simple-connection-zero calculation makes it nonzero. It includes both the connection slope and the source-rate clearing factor; replacing it by an unnamed Wronskian constant would lose its normalization.

\begin{theorem}[Five-factor survival criterion]\label{thm:survival}
At \(w=w_e(\rho)\), the physical Laurent coefficient is
\begin{equation}\label{eq:residue}
 \Res R_{q,b}(\,\cdot\,;a,x)
 =m_w(x)\frac{A_\rho M(w,s_b,a)O(s_b,w)}{D'_{q,b,e}}.
\end{equation}
For \(\re\rho>1/2\), \(R_{q,b}\) has no finite punctured-neighborhood limit at \(w_e\) if and only if
\begin{equation}\label{eq:five}
 \begin{gathered}
 \xi(5\rho/2-1)\ne0,\qquad\xi(\rho/2)\ne0,\qquad\xi(3\rho/2)\ne0,\\
 M(w_e,s_b,a)\ne0,\qquad O(s_b,w_e)\ne0.
 \end{gathered}
\end{equation}
\end{theorem}
\cert{A5}
\begin{proof}
At a fundamental connection zero, the two endpoint source integrals combine into the elementary nonzero mode \(m_w\) times a single moment. To see why the moment is independent of the interior splitting point, differentiate their sum: the derivative of the lower finite-part integral is the original integrand, and the derivative of the upper integral is its negative. The sum is constant and equals \eqref{eq:moment} at \(b=1/2\). The cleared Green numerator therefore factors as \(m_wMO\). Dividing by the simple denominator zero gives \eqref{eq:residue}. On \(x<0\), both bases in \(m_w\) are positive and nonzero; the denominator slope is nonzero. A simple-pole quotient has a nonremovable singularity exactly when its numerator at the pole is nonzero. Theorem~\ref{thm:angular} then yields \eqref{eq:five}.
\end{proof}

The word ``criterion'' is deliberate. No universal nonvanishing of the last two factors is asserted. The theorem allows \(a=0\), in which case the moment vanishes; a discussion of excitation must first choose a nonzero amplitude. A nonzero radial norm, evaluated next, is not a nonzero source moment.

\section{Radial regularization and Laurent extraction}

The mode-square finite part in the corpus is
\begin{equation}\label{eq:betanorm}
 N(w)=\frac{\Gamma(w)\Gamma(3/2)}{2\Gamma(w+3/2)},
 \qquad \re w>-2,\quad\im w\ne0.
\end{equation}
It is nonzero on that domain. This is the continued, horizon-subtracted mode square, not a positive Hilbert-space norm.
\cert{A8}
Its proof applies the endpoint-subtracted beta identity with parameters \(w,3/2\). The nonreal restrictions exclude all Gamma poles, and Gamma has no zeros, so numerator and denominator are nonzero. The analytic continuation and endpoint subtraction are part of the formal norm theorem.

There is also an actual integral--Laurent interchange for the horizon-mode temporal current. With \(w_0=w_e(\rho)\), let \(J(u,x)\) denote that current multiplied by its radius Jacobian, and let \(D(u)\) be its clearing factor. The theorem constructs analytic functions \(c_0(u),c_1(u)\) and a meromorphic function \(\Phi(u)\) such that, eventually for \(u\ne w_0\),
\begin{equation}\label{eq:radiallimit}
 \lim_{X\to-\infty}\left\{\int_X^0J(u,x)\dd x+
 \frac{c_0(u)z(X)^u/u+c_1(u)z(X)^{u+1}/(u+1)}{D(u)}\right\}
 =\Phi(u).
\end{equation}
Its Laurent coefficient of order \(-1\) is exactly the independently defined \code{physicalRadialResidueFinitePart}.
\cert{A9}

The proof first establishes a frequency-uniform second-order remainder bound for the cleared horizon density. Subtraction of the two displayed jets leaves an integrable Mellin remainder on the target frequency neighborhood. The complementary compact integral is analytic by domination. Factoring \(D(u)=(u-w_0)\kappa(u)\) with \(\kappa(w_0)\ne0\) identifies the residue from the cleared density at \(w_0\); uniqueness of its jet coefficients matches the prescribed counterterms. This argument justifies the interchange for this current. It gives meromorphy with a possible simple pole; it does not assert a nonzero residue when the source factor vanishes.

The branch selector in this result is \code{outer = false}. The corresponding normalizable-mode leg, all angular integration, and limits in time or the null horizon are not conclusions of \eqref{eq:radiallimit}. Iterating these operations requires its own uniform estimates and identification of common counterterms.

\section{Particular inverses and the literal quadratic metric}

\subsection{A finite Laurent particular solution}

Put \(f=r^2-1\), \(\lambda_2=(2s+1)(2s+2)\), and define
\[
 D_s=\partial_r+\frac{2s}{rf},\qquad
 \mathcal L_{s,\mu}=D_s^2+\frac{2r}{f}D_s+
                      \frac{\mu}{r^2f}-\frac{4s^2}{f^2}.
\]
The twist is essential: it differentiates the retained radial envelope. Set
\begin{align*}
 B&=\frac{s^2(s+1)^2(\mu-2s^2-6s-6)}{8(2-\mu)},\\
 M_s(\mu)&=-8s+2s\mu+s\mu^2-\tfrac52s^2-\tfrac94s^2\mu
 +\tfrac52s^2\mu^2+\tfrac{79}2s^3-\tfrac{17}2s^3\mu\\
 &\hspace{8mm}+54s^4-\tfrac{29}4s^4\mu+\tfrac{51}2s^5+\tfrac{11}2s^6,\\
 A&=\frac{-M_s(\mu)/(2-\mu)-6B}{\mu-\lambda_2}.
\end{align*}
Here \(M_s(\mu)\) is a polynomial, unrelated to the finite-part source moment \(M(w,s,a)\).

\begin{theorem}[Scalar--scalar Laurent inverse]\label{thm:laurent}
If \(\mu\ne2,\lambda_2\), then \(Y=r+A/r+B/r^3\) satisfies the original reduced scalar--scalar source equation
\[
 \mathcal L_{s,\mu}Y+\mathcal Q_{s,\mu}=0\quad(r>1),
\]
where the source has the exact decomposition
\[
 r^2f\mathcal Q_{s,\mu}=-2r^3-(\mu+2s-4s^2)r
 +\frac{M_s(\mu)}{(2-\mu)r}
 -\frac{[\mu-(2s+3)(2s+4)]B}{r^3}.
\]
At either physical parent rate, the associated sourced reconstruction data exist for every independent angular value \(\mu=-k^2\), \(k\in\R\).
\end{theorem}
\cert{B1}
\begin{proof}
Twice differentiating the Laurent polynomial, including the twist, gives
\[
 r^2f\mathcal L_{s,\mu}Y=2r^3+(\mu+2s-4s^2)r
 +\frac{(\mu-\lambda_2)A+6B}{r}
 +\frac{[\mu-(2s+3)(2s+4)]B}{r^3}.
\]
The definition of \(A\) cancels the middle source coefficient; the other coefficients cancel directly. The source decomposition itself follows by clearing denominators in the original source rows, with no removal of their trace-free term. The Laurent function is smooth on \(r>1\). For physical rates \(\re s<-1\), \(\im s\ne0\), and
\(\im\lambda_2=(8\re s+6)\im s\ne0\). Thus \(-k^2\ne\lambda_2\), and \(-k^2\ne2\); also \(s\ne0\). These are exactly the denominator and regularity requirements of the sourced reconstruction.
\end{proof}

At the fundamental harmonic \(\mu=(s+1)(s+2)=\rho(\rho-1)\), the first inverse-power coefficients simplify to
\begin{align*}
 A^{SS}_+(\rho)&=-\frac{(5\rho-7)(\rho+2)}{2(\rho+1)},&
 A^{SS}_-(\rho)&=\frac{(5\rho+2)(\rho-3)}{2(\rho-2)},\\
 A^{VV}_+(\rho)&=\tfrac23(\rho-2)(4\rho+3),&
 A^{VV}_-(\rho)&=\tfrac23(\rho+1)(4\rho-7).
\end{align*}
The last row concerns the separate vector--vector particular field.
\cert{B2}
These formulas determine particular coefficients. They do not, by themselves, prove retarded uniqueness or admission of the forced growing term \(r\) in the desired boundary class. A previous claimed obstruction to rational scalar--scalar reconstruction was retracted in the working notes after correction of untwisted derivatives and a doubled trace-free coefficient. The positive survivor is the explicit inverse above.

\subsection{The nonseparated scalar--scalar metric}

The independent-harmonic construction is not the endpoint of the source. Let
\[
 F=S_\rho^2,\qquad Q=\code{coordinateResponse}(q),\qquad
 R=\frac{2}{\lambda_2}\code{fullAngularField}(q,b).
\]
These are constructed smooth angular fields, and their equations are
\begin{equation}\label{eq:threeangular}
 \Delta Q=2Q-F,\qquad \Delta R=\lambda_2R+F.
\end{equation}
In particular, they are not arbitrary inverse-oracle arguments to the final theorem. Appendix~\ref{app:metric} gives the radial coefficient recipe. For \(Y_0=F,Y_2=Q,Y_D=R\), form
\begin{equation}\label{eq:metric}
 K=a^2e^{2s\tau}(1-r^{-2})^s
 \sum_{k\in\{0,2,D\}}Y_k
 \begin{pmatrix}
 F_{0k}&F_{1k}&0&0\\
 F_{1k}&F_{2k}&0&0\\
 0&0&2r^2H_k/y^2&0\\
 0&0&0&2r^2H_k/y^2
 \end{pmatrix}.
\end{equation}

\begin{theorem}[Literal sourced Einstein equation]\label{thm:metric}
For every Xi zero copy, either branch, every complex amplitude, and every exterior point \(r>1,y>0\), the twice differentiable field \eqref{eq:metric} satisfies all sixteen coordinate equations
\[
 D\mathcal E_{g_0}[K]+\mathcal Q^{SS}_{q,b,a}=0,
 \qquad \mathcal E(g)=G(g)-3g.
\]
Here \(\mathcal Q^{SS}_{q,b,a}\) is the actual scalar--scalar quadratic Einstein source, normalized as one half of the second variation of the original first-order scalar metric.
\end{theorem}
\cert{B3}
\begin{proof}[Proof structure]
Differentiate each of the three actual metric columns, retaining its envelope. The coordinate linearized operator is linear in these columns. Substitute \eqref{eq:threeangular} and the scalar equation for \(S_\rho\). Its angular terms reduce to \(F\), \(\Delta F\), its first derivatives, and the trace-free Hessians of \(F\) and \(Q\). The source has exactly the same tensor structures. The radial recipe in Appendix~\ref{app:metric} makes each coefficient of their sum zero. The corpus proves this coefficient cancellation before substituting the actual smooth fields, then identifies their two-jets with the original coordinate Einstein source. Its final theorem quantifies over both coordinate indices; no sector equation is left as an argument. The full component algebra and angular inverse constructions are in the accompanying dependency source snapshot.
\end{proof}

This closes the local, nonseparated scalar--scalar equation at the stated source rate. It does not construct a switched, retarded quadratic metric in every sector, the opposite-branch correction, or its complete integrated cubic action. Consequently the remaining action problem cannot be described as merely finding this scalar--scalar particular metric, but it also cannot be declared solved by its existence.

\section{The cubic coefficient retains the changing test weight}

Write \(L,B,C\) for the first, second, and third derivatives of \(\mathcal E\) at \(g_0\). Consider the coefficient of \(\epsilon^2\eta\) in a metric family with terms
\[
 g_0+\epsilon p+\eta m+\epsilon^2k_{++}
       +\epsilon\eta k_{+-}+\epsilon^2\eta\ell.
\]
Let \(W(g)\) be the metric-dependent contraction against a fixed test tensor, \(W_0=W(g_0)\), and \(W_1[h]=DW_{g_0}[h]\). For the original linearized solutions \(p,m\), the exact local paired coefficient is
\begin{align}\label{eq:cubic}
 &W_0\!:\!\left(L\ell+B(p,k_{+-})+B(m,k_{++})+\tfrac12C(p,p,m)\right)\nonumber\\
 &\quad+W_1[m]\!:\!\left(Lk_{++}+\tfrac12B(p,p)\right)
       +W_1[p]\!:\!\left(Lk_{+-}+B(p,m)\right).
\end{align}
The colon denotes the full coordinate contraction. The source theorem requires the displayed correction fields to be twice differentiable at the point. Their nonlinear field equations are not required for \eqref{eq:cubic}.
\cert{D2}
Taylor expansion proves the formula: the cubic tensor term has multiplicity three in the symmetric third derivative, hence coefficient \(3/3!=1/2\); the two products of a first weight variation and a quadratic residual give the last line. Terms involving a first-order residual vanish because the original first-order modes solve the linearized equation.

\textbf{Conditional simplification.} If both \(Lk_{++}+B(p,p)/2=0\) and \(Lk_{+-}+B(p,m)=0\) hold, the last line vanishes. This is the explicit premise of \code{mixedPairedCoefficient_of_quadratic_zero}, not an unconditional physical extinction theorem.
\cert{D3}
For the affine family \(k_{++}=k_{+-}=\ell=0\), the corpus instead gives an unconditional formula retaining \(W_0:C/2\), \(W_1[m]:B(p,p)/2\), and \(W_1[p]:B(p,m)\). Dropping the latter two is incorrect. Passing from this local density to an integrated action further requires a measure, domain, and boundary prescription.

\section{Angular current and the renormalized modular boundary}

\subsection{A bilinear source-contact identity}

For the literal scalar field, abbreviate
\begin{align*}
 E&=y^2(S_x^2+S_y^2),&
 C_0&=\frac{3\lambda-2}{16}S^2-\frac E4,\\
 B_0&=y^2(S_x(C_0)_x+S_y(C_0)_y),&
 G_0&=y^2((C_0)_x^2+(C_0)_y^2).
\end{align*}
The full source-contact calculation yields
\begin{align}\label{eq:contact}
 &(\lambda-2)(\lambda-6)(3\lambda-2)S^4\nonumber\\
 &=16(\lambda-2)(\lambda-6)S^2C_0
 +64(\lambda-2)(2SB_0+C_0E)\nonumber\\
 &\hspace{8mm}-128E\Delta C_0-512G_0.
\end{align}
Its prefactor is nonzero at every Xi zero: \(\re\lambda=\sigma(\sigma-1)-t^2<0\), whereas its three possible roots are positive real numbers. The identity follows by applying the angular equation and product rule to the displayed contact and retaining all gradient terms.
\cert{D4}
All products in \eqref{eq:contact} are complex bilinear. Thus \(E\) and \(G_0\) need not be nonnegative. This rigidity identity cannot supply a positive integrated energy without further structure.

\subsection{The curved modular domain}

For \(Y\ge1\), use the actual truncated modular domain
\[
 \mathcal D_Y=\{(x,y):|x|\le1/2,\ x^2+y^2\ge1,\ y\le Y,\ y>0\}.
\]
At fixed null parameters \(p,v\), let \(F\) be the original scalar potential and let \(T=S_u\) be the testing Eisenstein scalar, \(\Delta T=\mu T\), \(\mu=u^2-1/4\). Define a coordinate flux
\[
 j=(\Delta F)\nabla T-\mu T\nabla F.
\]
Its Euclidean divergence is the source's coexact-vector defect. The measure in the next formula is \(\dd x\dd y\), not \(y^{-2}\dd x\dd y\).

\begin{theorem}[Original modular Green finite part]\label{thm:modular}
For actual source and reader Xi copies, either source branch and any amplitude,
\[
 \int_{\mathcal D_Y}\operatorname{div}j\dd x\dd y
       =\int_{-1/2}^{1/2}j_y(x,Y)\dd x.
\]
For fixed \(p,v\ge0\), subtract from this top flux the corpus's grouped nonconstant cusp-exponent family. Its limit at \(Y=e^k\), \(k\to\infty\), equals the coefficient of exponent zero. That coefficient, and hence this prescribed finite part, is zero when reader copy equals source copy or its functional-equation reflection.
\end{theorem}
\cert{B4}
\begin{proof}[Proof structure]
The scalar and coexact-vector identities identify the integrand with the coordinate divergence. Modular translation cancels the vertical sides. Modular inversion identifies the two pieces of the curved lower boundary with opposite flux; replacing the domain by a rectangle would lose this step. Stokes' formula leaves only the top edge. The actual Fourier--Bessel tail has the decay needed for its contribution to the subtracted limit to vanish. Grouping equal exponents before subtraction leaves the coefficient of exponent zero, including coincident-exponent cases. The diagonal and reflected zero-coefficient identities then give the final assertion. The formal statements apply to the original source, including onset and both Weyl terms.
\end{proof}

The finite part is a specified subtraction scheme, not the assertion that the unrenormalized integral converges to zero. Nor does a zero angular finite part at each fixed pair of parameters license exchanging that limit with an improper time integral or horizon limit.

\subsection{The actual original metric current}

There is also an exact local identification for the original physical angular current. In ingoing coordinates \((p,R,x,y)\), put \(U=RZ\), where \(Z\) is the original ingoing potential, and let \(W\) be the original left scalar with its physical radial normalization. For angular index \(j\), define
\[
 \mathcal B_j(F,G)=(\Delta F)\partial_jG-(\Delta G)\partial_jF.
\]
Then the original metric current is
\begin{equation}\label{eq:originalcurrent}
 J^j=\frac1{R^2}\left[(R^2-1)\mathcal B_j(U_R,W_R)
                 +\mathcal B_j(U_R,W_p)+\mathcal B_j(U_p,W_R)\right].
\end{equation}
The source and reader copies may be independent; both original testing sides and both radial normalizations are included. The proven domain is \(p>0,1<R<3,y>0\).
\cert{B5}
The proof commutes the actual mixed derivatives, identifies the causal stream as the coexact derivative of \(RZ\), and substitutes in the four-component Einstein current. The source and left-field smoothness are supplied by their constructions. In particular, differentiating \(Z\) while discarding the prefactor \(R\) does not reproduce this current.

Equation~\eqref{eq:originalcurrent} is a more faithful starting point than a projected current with angular components set to zero. Its derivative-weighted Green transfers still require uniform parameter-dependent boundary estimates before Theorem~\ref{thm:modular} can be used inside the complete spacetime action. That transfer is an open problem here.

\section{An exact rational detector and its limits}

The elementary plateau calculations provide a useful example of what an unconditional detector does, and what it does not do. Let \(K(s)\) be the value of \code{skewPlateauMu} at angular parameter zero. Set \(v=s+3/2\) and
\begin{align*}
 P(u)&=-\tfrac{74655}{2048}+\tfrac{58863}{256}u-\tfrac{19053}{32}u^2
 +\tfrac{36367}{32}u^3-\tfrac{3151}{8}u^4-46u^5+36u^6-8u^7,\\
 Q(u)&=4-20u+16u^2=(4u-4)(4u-1).
\end{align*}
Away from the four elementary denominator zeros,
\begin{equation}\label{eq:kernel}
 K(s)=\frac{vP(v^2)}{512Q(v^2)}.
\end{equation}
The denominator zeros are \(s=-1/2,-1,-2,-5/2\); none lies at positive imaginary height.

\begin{theorem}[Constant-mode tilt detection]\label{thm:tilt}
For \(0\le\sigma\le1\) and \(t\ge92\),
\[
 \re K(\sigma-2+it)=0\quad\Longleftrightarrow\quad\sigma=1/2.
\]
More precisely, there is an explicit real polynomial \(C(\sigma,t)>0\) on this domain such that
\begin{equation}\label{eq:tilt}
 2\re K(\sigma-2+it)\,|512Q((\sigma-1/2+it)^2)|^2
      =(2\sigma-1)C(\sigma,t).
\end{equation}
\end{theorem}
\cert{C3}
\begin{proof}
Define the polynomial \(H\) through the symmetric difference identity
\begin{align*}
 &vP(v^2)512Q(y^2)-512Q(v^2)yP(y^2)\\
 &\hspace{18mm}=(v-y)512H((v+y)^2,vy).
\end{align*}
The corpus gives its explicit coefficients and verifies this polynomial identity. This definition also permits direct calculation of \(H\) by polynomial division and symmetric reduction. Equation~\eqref{eq:kernel} is odd under \(s\mapsto-3-s\) and commutes with conjugation. Choose \(v=\sigma-1/2+it\) and \(y=-\overline v\). Then
\[
 (v+y)^2=-4t^2,\quad vy=-((\sigma-1/2)^2+t^2),\quad v-y=2\sigma-1.
\]
Taking the difference of the two rational kernels gives \eqref{eq:tilt} with
\[
 C(\sigma,t)=512H(-4t^2,-((\sigma-1/2)^2+t^2)).
\]
For completeness, the elementary bound used for positivity is
\[
 C(\sigma,t)\ge720896t^{18}-6067244032t^{16}
 \quad(0\le\sigma\le1,\ t\ge1).
\]
Expand \(C\) in \(\sigma,t\), discard nonnegative monomials, and bound the remaining monomials by \(t^{16}\). The four coefficient groups in the formal certificate have absolute negative sums
\[
 2244870144,\quad1077216256,\quad1209179136,\quad1535978496.
\]
Their sum is the displayed constant. Since \(720896\cdot92^2>6067244032\), the bound is positive at \(t\ge92\). The squared denominator is positive, so \eqref{eq:tilt} proves the equivalence.
\end{proof}

This theorem holds for every complex parameter in its domain, regardless of whether it is a Xi zero. It supplies a rational test for the real coordinate. It supplies no reason that its value at an arithmetic zero must vanish, and it makes no claim below height 92.

\subsection{Amplitude dependence and spectral assembly}

The scalar, mixed, and vector constant-mode leading coefficients have the exact factorization
\begin{equation}\label{eq:square}
 8x^2-192xy+1152y^2=8(x-12y)^2.
\end{equation}
Here \(x,y\) stand for the squared scalar and vector amplitudes in the leading quartic form.
\cert{C4}
For real \(x,y\), this is nonnegative, but it vanishes on \(x=12y\). Consequently it cannot establish the notes' assertion of an automatic strict sign in every reading. For complex amplitudes, even the word nonnegative requires a separate conjugation identification. A helicity substitution in a bilinear polynomial produces a fourth power of the amplitude; replacing that with an absolute fourth power also requires such an identification.

The notes propose that a physical coefficient has a dominant constant-mode contribution of schematic form
\[
 K_{\rm phys}(\rho)=c_0(\rho)^2K(\rho-2)+\mathcal R(\rho).
\]
\textbf{Unproved assembly claim.} This article does not assert this decomposition for the complete original action, nonvanishing of its normalizing coefficient, or the proposed relative error \(O(t^{-3/2})\). The isolated harmonic estimates do not supply the regularized spectral identity, the measure, the mixed-sector amplitudes, and the boundary interchange needed here.

There is a second issue even if an asymptotic estimate is established. The leading real part in \eqref{eq:tilt} tends to zero with \(2\sigma-1\). A small additive error bound alone cannot justify a uniform sign arbitrarily close to the line. A valid transfer needs a matching tilt factor in the real error, or a quantitative bound relative to that vanishing leading real part. A magnitude estimate relative to the full complex kernel is insufficient. The same concern applies when a cutoff is chosen separately for each zero and then promoted to a single universal choice.

\section{The exact boundary of the original source probe}

In the corpus's Peierls carrier, write \(s_+(q)\) and \(s_-(q)\) for the two unit source vectors in this paragraph only; these vectors are unrelated to the scalar rates \(s_\pm(\rho)\). Let \(\beta\) be its bilinear pairing. The complementary testing vector obeys
\[
 s_-(q)=0\quad\Longleftrightarrow\quad\re\rho=1/2.
\]
For arbitrary returned vectors \(r(q)\), define the measured profile \(T_r(q)=\beta(r(q),s_-(q))\).

\begin{theorem}[Exact attainable profiles]\label{thm:range}
A function \(T:\mathrm{XiZeroCopy}\to\C\) equals \(T_r\) for some returned-vector family if and only if it vanishes on all critical-line copies. Every such profile is realized by a weighted complex-linear skew propagator applied to the original positive source.
\end{theorem}
\cert{C1}
\begin{proof}
Necessity follows from pairing against zero. For sufficiency, let
\(h(q)=|A_{\rm pol}(q)|^2\), where \(A_{\rm pol}\) is the unit-amplitude value of \code{d4BMXiSecondHarmonicPolarizedSourceAmplitude}.
The source identities give \(h(q)\ne0\) exactly off the critical line, and the weighted propagator has profile \(w(q)h(q)\). Choose \(w(q)=T(q)/h(q)\) off the line, and zero on it. This produces \(T\) everywhere. No regularity or linearity of an arbitrary competing returned-vector family was used.
\end{proof}

This is both a no-go and a positive classification. Nonlinear propagation cannot recover information already annihilated by this particular testing leg. Conversely, the canonical construction realizes every compatible target. Because it uses that target in its defining weight, it is not an independently derived physical response law. The theorem excludes neither a different testing field nor a retained derivative of a vanishing source; those possibilities need their own construction and physical identification.

\section{Receivers for the Riemann hypothesis}

The phase-sensitive chiral Ward construction is particularly transparent. With a nonzero source amplitude \(a\), its reduced projection has the factorization
\[
 \mathcal W_a(q)=m_a(q)(1-2\re\rho),\qquad m_a(q)\ne0.
\]
In the source normalization,
\[
 m_a(q)=i\im\rho\,
   \code{d4BMGaugeFixedSecondHarmonicGreen}(q)
   C_{\mathrm{sd}}(a,z_q).
\]
Here \(C_{\mathrm{sd}}\) denotes \code{d4ThetaHawkingSelfDualQuadraticCarrierCoefficient}.
Dividing by this multiplier gives a real calibrated value \(1-2\re\rho\). Therefore the corpus proves
\[
 \bigl(\forall q,\ \mathcal W_1(q)=0\bigr)\quad\Longleftrightarrow\quad\RH.
\]
\cert{C2}
The proof is cancellation of the nonzero multiplier followed by the zero-copy formulation of \(\RH\). It is a premise-free equivalence, but neither side is supplied as a premise-free fact. Identifying \(\mathcal W\) with the complete physical mixed action is not a consequence of defining this reduced projection.

The static front has analogous equivalences for tilted-resonance centralizer rigidity and KMS Cayley scattering mixing.
\cert{C7}\cert{C8}
Their role is to translate a universal property of a zero-labelled construction into zero location. Abstract graph completion, an operator realization, or a local near-action law does not establish that universal property. These alternatives remain separate possible endpoints until a source-derived dynamical or static principle supplies one of them.

\section{Arithmetic nonvanishing and the unresolved frontier}

\subsection{What affine collision control achieves}

The affine arguments in \eqref{eq:collisions} introduce a separate arithmetic obstacle. The corpus's half-descent sends \(\rho\) to \(1-\rho/2\) when the corresponding half-factor is a zero. Iteration approaches the real fixed point \(2/3\); isolation of zeros and nonvanishing there prevent indefinite descent. This can produce a strict-right zero whose half-factor is nonzero. It does not assert that the other two factors remain nonzero at the chosen descendant.
\cert{A10}

For the five-half factor, the maps used are \(\rho\mapsto5\rho/2-1\) and the reflected ascent \(\rho\mapsto2-(5/2)\overline\rho\). The precise source statements are as follows.

\begin{proposition}[Five-half collision alternatives]
Suppose \(\xi(\rho)=0\), \(\re\rho>1/2\), and \(\im\rho\ne0\). If \(\re\rho\notin[8/15,2/3]\), there exists a strict-right nonreal zero \(\rho'\) with both
\(\xi(5\rho'/2-1)\ne0\) and \(\xi(3\rho'/2)\ne0\).
Without the band exclusion, the formal orbit argument gives three alternatives: a strict-right zero with nonzero five-half factor exists; an iterate has real part \(3/5\); or every iterate remains a zero with real part in \([8/15,2/3]\).
\end{proposition}
\cert{C5}
The proof follows the affine collision orbit until it escapes the allowed zero strip, unless it meets the specified band or line. The fixed real coordinates explain why the escape argument does not close the trapped alternatives. The source supplies the precise orbit and induction, rather than an assertion that all three affine collision conditions are independent.

Even the declaration named \code{physicalSelfSourceCoefficient_ne_zero_of_four_fifths_le} requires both \(\re\rho\ge4/5\) and the argument \(\xi(\rho/2)\ne0\). Its name must not be read as unconditional nonvanishing in that substrip.
\cert{C6}

\subsection{Explicit open problems}

\begin{problem}[Simultaneous excitation]\label{prob:excite}
Starting from a hypothetical off-line zero, produce a zero and physical branch for which all five conditions in \eqref{eq:five} hold simultaneously, with a nonzero amplitude and one admitted time convention. Resolve the trapped affine collision alternatives or supply a different source functional. Prove nonvanishing of the actual moment \eqref{eq:moment}; a different radial moment or the norm \eqref{eq:betanorm} is insufficient.
\end{problem}

\begin{problem}[Onset transport]
Re-derive the full causal pipeline with \(\chi(u/L)\), including the source, Green response, counterterms, and clock identity, at a duration furnished by Theorem~\ref{thm:onset}. Alternatively prove nonvanishing at the already implemented unit duration. Merely substituting \(L\) in the scalar pulse transform does not complete this task.
\end{problem}

\begin{problem}[Complete nonlinear action]\label{prob:action}
Construct all required quadratic corrections with the original source, gauge, boundary condition, and retarded onset. In particular identify the opposite-branch correction used by \eqref{eq:cubic}. Establish the full integrated coefficient with both changing-weight terms accounted for. Relate the particular metric in Theorem~\ref{thm:metric} to that causal solution by an actual admission and uniqueness theorem.
\end{problem}

\begin{problem}[Compatible limiting operations]
Extend the radial Laurent extraction to the required normalizable testing leg. Prove the derivative-weighted cusp transfers for \eqref{eq:originalcurrent}, the horizon and boundary trace estimates, and a common subtraction prescription allowing the angular, radial, temporal, and Laurent operations in the physical action. Fixed-parameter finite parts alone do not establish these interchanges.
\end{problem}

\begin{problem}[A faithful physical left reader]
Construct a left testing object that retains the nonzero first live spin-two jet, or prove that the desired physical coefficient needs only the information preserved by the original probe. Specify its differential equation, domain, normalization, and causal pairing. The exact range obstruction in Theorem~\ref{thm:range} must be respected.
\end{problem}

\begin{problem}[Spectral and helicity matching]
Identify the physical scalar/vector amplitudes and the exact integrated kernel before using per-harmonic sign laws. If the constant-mode approach is retained, prove the regularized decomposition, nonzero normalization, tail bound, and a tilt-relative real error estimate, including degenerate amplitude ratios. Match the metric harmonic parameters to the original optical testing frequencies rather than equating them by notation.
\end{problem}

\begin{problem}[Extinction from the source]\label{prob:producer}
Prove a source-derived law that makes a faithful complete physical diagnostic vanish, without assuming \(\RH\), its equivalent mixing or Ward condition, or the desired nonvanishing contradiction. For the static alternative, derive the requisite phase-sensitive KMS or centralizer property from the actual source. The equivalences in the preceding section are receivers for such a law.
\end{problem}

The proposed reduction would combine a faithful nonzero response at a hypothetical off-line zero with an independent law forbidding that response. Neither the forbidden response nor the law is supplied by naming a kernel ``physical'' or a condition ``extinction.'' The problems above specify where the reduction currently lacks that content.

The central risk is conceptual as well as analytic. Completing the existing inverse, regularity, and boundary constructions would give a better-defined action, but would not by itself derive a selective vanishing law. Conversely, a vanishing theorem for an abstract or projected expression is irrelevant to the route unless it is identified with the correct physical diagnostic. No claim that these two requirements can be met simultaneously follows from the current corpus audit. This is an unresolved viability question, not a theorem that the route is impossible.

\section{Coverage and corrections to the working ledger}

Table~\ref{tab:ledger} addresses every numbered ledger item, including 33a and 33b. ``Supporting development'' describes the ledger's subsidiary analytic or algebraic package; it does not promote every sentence of that row to a theorem of this article. Exact declarations for results asserted here are in Appendix~\ref{app:cert}. The supplementary ledger review retains the complete input rows, source-anchor matches, and these scope decisions. In particular, a word such as ``FULL'' in the input is not an axiom or dependency audit.

\small
\begin{longtable}{@{}p{.07\linewidth}p{.28\linewidth}p{.59\linewidth}@{}}
\caption{Item-by-item Route HB assessment.}\label{tab:ledger}\\
\toprule Item & Topic & Audited interpretation and remaining restriction\\\midrule
\endfirsthead
\toprule Item & Topic & Audited interpretation and remaining restriction\\\midrule
\endhead
\bottomrule\endfoot
1 & Literal theta source & Retain both Weyl terms and tail; scalar equation is a source result. Joint smoothness is local, not cusp integrability.\\
2 & Xi occurrence & Optical incidence, including multiplicities and shadow, is Proposition~\ref{thm:incidence}; occurrence does not force a pole to disappear.\\
3 & Excitation/Jost residues & Supporting Hardy/Jost development concerns its specified source and boundary problem; it does not discharge the five physical factors.\\
4 & Spin-two source & Entire factorization and strip unit are Theorem~\ref{thm:amputation}; retain the external Xi square and distinguish the first live jet.\\
5 & Angular source value & Exact formula and collision criterion, Theorem~\ref{thm:angular}; nonvanishing remains separate.\\
6 & Metric jets & Supporting coordinate linearization and multilinear variations; actual first-order equations underlie \eqref{eq:cubic}.\\
7 & Mixed quadratic block & A sector solution does not supply every quadratic slot in the mixed action.\\
8 & Cubic coefficient & Exact local formula \eqref{eq:cubic}; deletion of changing-weight terms is conditional on both quadratic equations.\\
9 & Polarization obstruction & A scoped loss of tensor data, not impossibility of every scalar functional. Retain the actual source tensor.\\
10 & Angular inverse & Constructed inverses now enter \eqref{eq:threeangular}; operator existence does not imply causal extinction.\\
11 & Time source & Supporting onset and Laplace admission at the implemented convention; universal-duration transport is open.\\
12 & Native inverse & Supporting Legendre/Volterra realization; source normalization and each principal channel matter.\\
13 & Energy/Friedrichs & Supporting residual-space admission; it is not positivity of a non-Hermitian QNM action.\\
14 & Native/physical equality & Supporting almost-everywhere identification; continued poles still require the survival criterion.\\
15 & Green band & Equality is on its stated Laplace strip; continuation and nonvanishing are additional steps.\\
16 & QNM occurrence & Fixed-domain connection zeros do not by themselves imply excited poles of a forced response.\\
17 & Laurent factorization & Theorem~\ref{thm:survival}; all five factors are retained.\\
18 & Action/current & Local identities precede the integrated boundary and limiting problems.\\
19 & Projected action & Vanishing angular components belong to the projected pair; use \eqref{eq:originalcurrent} for the original current.\\
20 & Temporal residue & Local continuation and factorization; no integrated nonzero action follows automatically.\\
21 & Radial finite part & Nonzero continued mode square \eqref{eq:betanorm}; it is not a nonzero source moment.\\
22 & Incoming expansion & Uniform endpoint jets support the scoped Laurent theorem, not every iterated limit.\\
23 & Contact horizon factor & Cancellation of an extra pole is a local contact result, not extinction of the surviving response.\\
24 & Pulse/history & Zero contact moment does not erase retarded history. A proposed clock proof must retain it.\\
25 & Late clock & The terminal exponential cannot be omitted from the two-source expression. No general extinction inferred.\\
26 & Source translation & The fixed-clock contact moment obstructs naive switched-source covariance even on the line.\\
27 & Incoming moments & Supporting component and balance identities; identify the same fields and clock before integrating.\\
28 & Angular propagation & Supporting circle/Goursat kernel; transfer to the original modular domain is a further identification.\\
29 & Contact rigidity & Exact bilinear identity \eqref{eq:contact}; no positivity assertion.\\
30 & Spin-two balance & Paper-level proposal in the ledger. Not asserted as proved in this article.\\
31 & Laurent interchange & Horizon-mode leg only. Correct ``simple pole'' to ``possible simple pole'' unless residue is nonzero.\\
32 & QNM boundary germ & Scoped horizon phase and bracket results do not discharge every differentiated or cusp boundary face.\\
33a & Onset & One common positive duration exists; fixed unit nonvanishing and full duration transport remain open.\\
33b & Source moment & Nonvanishing of \eqref{eq:moment} remains unproved here; a different moment cannot replace it.\\
34 & Affine collisions & Partial escape and half-descent results; simultaneous control and trapped cases remain.\\
35 & Quadratic corrections & Update ledger: literal nonseparated scalar--scalar equation now exists. Retarded admission and opposite-branch action slot remain. Retract the old rational-inverse no-go.\\
36 & Angular/cusp integral & Update ledger: original curved-domain finite part and original angular-current identity exist. Their full spacetime transfer is still open.\\
37 & Complete coefficient & Particular harmonic formulas and rational tilt are results; full spectral dominance, normalized physical identification, and extinction are not established.\\
38 & Left reader & Exact source-probe image is Theorem~\ref{thm:range}; no arbitrary returned vector repairs a zero testing leg.\\
39 & AIME match & The proposed physical sign/period identity still needs amplitude, reader, spectral, and boundary matching. No replacement of bilinear powers by modulus squares.\\
40 & Extinction producer & No unconditional producer was established in this audit. Scoped no-gos cannot justify a universal assertion that all scalar or finite-energy mechanisms are impossible.\\
41 & Extinction/RH & Equivalences are receivers. Their universal hypotheses are not conclusions.\\
42 & Static completion & Supporting graph/operator constructions; completion is not the required selective source law.\\
43 & Static receivers & Centralizer and mixing equivalences do not independently establish their premises.\\
44 & Static producer & Phase-sensitive source law remains open; generic positivity alone does not supply it.\\
\end{longtable}
\normalsize

\section{Relation to existing analysis and scope of the contribution}

The linear topological-black-hole literature \cite{aros,ki,bm} motivates the separated operator, but does not prove a nonlinear arithmetic selection law. A comparison must preserve the perturbation sector and boundary condition: a closed frequency formula in a linear problem cannot be imported as an equality of complete forced actions.

Bilinear QNM products require a similar distinction. Green, Hollands, Sberna, Toomani, and Zimmerman construct a conserved bilinear product and renormalized mode relations for Kerr \cite{green}. The 2026 preprint of Minucci, Macedo, Pantelidou, and Sberna studies scalar perturbations of the Schwarzschild family on hyperboloidal foliations; finite mode functions still yield divergent naive bilinear integrands, motivating regularization and flux terms \cite{minucci}. These constructions explain why reflection, pairing, and boundary treatment deserve separate proofs. They are not formal certificates for the theta-cusp action considered here.

Lestingi, Sberna, and Green's 2026 perturbation framework explicitly retains continuous-spectrum contributions alongside QNMs and exhibits divergent mode sums in examples \cite{lestingi}. It therefore cannot be cited as an unconditional justification for a discrete QNM expansion of the HB nonlinear source. On the automorphic side, Zagier's regularization concerns non-rapidly-decaying functions \cite{zagier}. The 2026 work of Dobrowolski, Jana, and Nunes gives a regularized distributional decomposition for a squared magnitude of a degenerate Eisenstein series on \(\mathrm{GL}_n\), with a separate moment asymptotic conditional on generalized Ramanujan \cite{djn}. Its object is not the present bilinear square \(S_\rho^2\), and no transfer of its conclusions to the full original HB current is asserted.

The results presented here therefore constitute a structured formal-source exposition and a precise research frontier. The rational inverse, literal sourced metric, residue criterion, and exact detector can be assessed independently of the proposed RH application. A claim of a new RH proof would require solving Problem~\ref{prob:producer} together with the identifications and excitation obligations above. The present manuscript makes no such claim.

\appendix
\section{Source concordance and reproducibility}\label{app:cert}

The supplementary package fixes the source rather than relying on a moving working tree. The repository HEAD recorded at the start of the audit is
\begin{center}\path{6ba671fa1ddf6c1f519445b4a0172e2abbfcac1c}.\end{center}
The working tree contained additional changes, so this commit alone does not identify the results. The archived source hashes are authoritative for this manuscript's source references. The toolchain file specifies Lean 4.32.2; the included Lake manifest pins external packages. The local import closure is archived, while external dependency source is obtained from those pins.

The file \path{audit/declarations.json} supplies full declaration types, namespaces, source lines, and file hashes. It contains 67 selected declarations in 28 groups. The unexecuted \path{audit/Verify.lean} lists the corresponding checks and axiom queries. A future kernel run must use the supplied source snapshot and pinned dependencies and must inspect both the complete types and transitive axioms. No new Lean checking was completed during this preparation, and the interrupted log must not be used as evidence that the final declarations passed.

The global source lint run also did not pass: it reported two unreachable modules as well as size warnings. That structural failure is not a discovered proof of a false mathematical statement, but it precludes reporting a clean whole-project gate. Source searches and formula-transcription checks are useful cross-checks; neither substitutes for kernel validation.

Below, a namespace is followed by short declaration names. The name of a module is relative to \path{formal/GaloisForLFunctions/RHTower/}. Conditions in the main text remain part of the statement even when the short name omits them. Groups D2--D3 include explicitly conditional formulas. The ``source group'' labels are provenance references, not newly executed certificates.

\subsection*{Group D5}\phantomsection\label{cert:D5}
\noindent Module: \path{ThetaCuspHawkingBMAngularScalarEquation.lean}.\\
Namespace: \path{RHTower.D4BMAngularSeries}.
\begin{itemize}
\item \path{scalarCusp_hyperbolicLaplacian}.
\end{itemize}

\subsection*{Group A10}\phantomsection\label{cert:A10}
\noindent Module: \path{XiAffineCollisionControl.lean}.\\
Namespace: \path{RHTower}.
\begin{itemize}
\item \path{exists_xi_zero_half_ne_zero}.
\end{itemize}

\subsection*{Group A1}\phantomsection\label{cert:A1}
\noindent Module: \path{ThetaCuspHawkingTopologicalBlackHoleQNMIncidence.lean}.\\
Namespace: \path{RHTower}.
\begin{itemize}
\item \path{thetaCuspHawkingXiDivisor_topologicalBlackHoleQNM_incidence}.
\item \path{thetaCuspHawkingTopologicalBlackHoleSeparatedOpticalEnergy_shadow}.
\item \path{thetaCuspHawkingTopologicalBlackHoleBranchImaginarySplitting_eq}.
\end{itemize}

\subsection*{Group A2}\phantomsection\label{cert:A2}
\noindent Module: \path{ThetaCuspHawkingBMCanonicalSpinTwoAmputation.lean}.\\
Namespace: \path{RHTower}.
\begin{itemize}
\item \path{d4BMCanonicalSpinTwoAmputatedSource_eq}.
\item \path{d4BMCompletedSpinTwoVertex_eq_XiC_sq_mul_canonicalSource}.
\item \path{d4BMCanonicalSpinTwoAmputation_existsUnique}.
\item \path{d4BMCanonicalSpinTwoAmputatedSource_eq_zero_iff}.
\end{itemize}

\subsection*{Group A3}\phantomsection\label{cert:A3}
\noindent Module: \path{ThetaCuspHawkingBMPhysicalAngularSourceValue.lean}.\\
Namespace: \path{RHTower.D4BMJointSource}.
\begin{itemize}
\item \path{physicalSelfSourceCoefficient_fourLambda}.
\end{itemize}

\subsection*{Group A4}\phantomsection\label{cert:A4}
\noindent Module: \path{ThetaCuspHawkingBMPhysicalAngularSourceXi.lean}.\\
Namespace: \path{RHTower.D4BMJointSource}.
\begin{itemize}
\item \path{physicalSelfSourceCoefficient_zero_iff_xi}.
\end{itemize}

\subsection*{Group A5}\phantomsection\label{cert:A5}
\noindent Module: \path{ThetaCuspHawkingBMPhysicalFundamentalMoment.lean}.\\
Namespace: \path{RHTower.D4BMJointSource}.
\begin{itemize}
\item \path{physicalSelfResidueCoefficient_fundamental_value}.
\item \path{continuedPhysicalSelfResponse_nonremovable_iff_elementary}.
\end{itemize}

\subsection*{Group A6}\phantomsection\label{cert:A6}
\noindent Module: \path{ThetaCuspHawkingBMClearedOnsetPulseTransform.lean}.\\
Namespace: \path{RHTower.D4BMJointSource}.
\begin{itemize}
\item \path{clearedOnset_eq_onsetPulseTransform}.
\item \path{clearedOnset_source_rate}.
\item \path{clearedOnsetDuration_one}.
\item \path{clearedOnsetDuration_zero}.
\item \path{differentiable_clearedOnsetDuration}.
\end{itemize}

\subsection*{Group A7}\phantomsection\label{cert:A7}
\noindent Module: \path{ThetaCuspHawkingBMUniversalOnsetDuration.lean}.\\
Namespace: \path{RHTower.D4BMJointSource}.
\begin{itemize}
\item \path{exists_universal_onset_duration}.
\end{itemize}

\subsection*{Group A8}\phantomsection\label{cert:A8}
\noindent Module: \path{ThetaCuspHawkingBMPhysicalRadialNorm.lean}.\\
Namespace: \path{RHTower.D4BMFundamentalNorm}.
\begin{itemize}
\item \path{normFinitePart_eq_gamma}.
\item \path{normFinitePart_ne_zero}.
\end{itemize}

\subsection*{Group A9}\phantomsection\label{cert:A9}
\noindent Module: \path{ThetaCuspHawkingBMTemporalCurrentRadialLaurent.lean}.\\
Namespace: \path{RHTower.D4BMJointSource}.
\begin{itemize}
\item \path{continuedPhysicalTimeAction_horizon_radial_laurent}.
\end{itemize}

\subsection*{Group B1}\phantomsection\label{cert:B1}
\noindent Module: \path{ThetaCuspHawkingBMScalarStreamLaurentMaster.lean}.\\
Namespace: \path{RHTower.D4BMScalarStream}.
\begin{itemize}
\item \path{ssMasterSource_laurent}.
\item \path{laurentMasterA_resolvents}.
\item \path{laurentMasterB_resolvent}.
\item \path{laurentMasterField_operator}.
\item \path{ssLaurentMaster_equation}.
\item \path{physical_ssLaurentMaster_sourcedData}.
\end{itemize}

\subsection*{Group B2}\phantomsection\label{cert:B2}
\noindent Module: \path{ThetaCuspHawkingBMLaurentPhysicalRadiation.lean}.\\
Namespace: \path{RHTower.D4BMScalarStream}.
\begin{itemize}
\item \path{vvLaurentMasterA_physical_plus}.
\item \path{vvLaurentMasterA_physical_minus}.
\item \path{laurentMasterA_physical_plus}.
\item \path{laurentMasterA_physical_minus}.
\end{itemize}

\subsection*{Group B3}\phantomsection\label{cert:B3}
\noindent Module: \path{ThetaCuspHawkingBMThreeFieldMetric.lean}.\\
Namespace: \path{RHTower.D4BMThreeFieldMetric}.
\begin{itemize}
\item \path{literalProfile_equations}.
\item \path{literalMetric_einstein}.
\end{itemize}

\subsection*{Group B4}\phantomsection\label{cert:B4}
\noindent Module: \path{ThetaCuspHawkingBMOriginalModularGreen.lean}.\\
Namespace: \path{RHTower.D4BMOriginalModularGreen}.
\begin{itemize}
\item \path{original_modular_area}.
\item \path{original_modular_finitePart}.
\item \path{diagonal_modular_finitePart}.
\item \path{reflected_modular_finitePart}.
\end{itemize}

\subsection*{Group B5}\phantomsection\label{cert:B5}
\noindent Module: \path{ThetaCuspHawkingBMOriginalAngularCurrent.lean}.\\
Namespace: \path{RHTower.D4BMOriginalAngularCurrent}.
\begin{itemize}
\item \path{original_physical_current}.
\item \path{original_physical_metric_current}.
\end{itemize}

\subsection*{Group C1}\phantomsection\label{cert:C1}
\noindent Module: \path{ThetaCuspHawkingBMSourceProbeRangeClassification.lean}.\\
Namespace: \path{RHTower}.
\begin{itemize}
\item \path{hbXiPeierlsSourceProbeRealizingWeight_profile}.
\item \path{exists_hbXiPeierlsSourceProbeProfile_iff}.
\end{itemize}

\subsection*{Group C2}\phantomsection\label{cert:C2}
\noindent Module: \path{ThetaCuspHawkingBMPhaseSensitiveChiralWardAmplitude.lean}.\\
Namespace: \path{RHTower}.
\begin{itemize}
\item \path{hbXiPhaseSensitiveChiralWardMultiplier_ne_zero}.
\item \path{d4BMMixedHorizonOpticalSourceProjection_eq_chiralMultiplier_mul_deningerCharge}.
\item \path{hbXiCalibratedPhaseSensitiveChiralWardAmplitude_re_eq}.
\item \path{hbXiCalibratedPhaseSensitiveChiralWardAmplitude_im_eq_zero}.
\item \path{riemannHypothesis_iff_hbXiPhaseSensitiveChiralWardAmplitudeCancellation}.
\end{itemize}

\subsection*{Group C3}\phantomsection\label{cert:C3}
\noindent Module: \path{ThetaCuspHawkingBMEvenSectorSkewConstantTilt.lean}.\\
Namespace: \path{RHTower.D4BMScalarStream}.
\begin{itemize}
\item \path{skewConstantKernel_eq}.
\item \path{skewConstantKernel_polynomial_identity}.
\item \path{skewConstantKernel_re_tilt}.
\item \path{skewConstantTiltCofactor_pos}.
\item \path{skewConstantKernel_re_eq_zero_iff}.
\end{itemize}

\subsection*{Group C4}\phantomsection\label{cert:C4}
\noindent Module: \path{ThetaCuspHawkingBMEvenSectorSkewConstantSquare.lean}.\\
Namespace: \path{RHTower.D4BMScalarStream}.
\begin{itemize}
\item \path{skewConstantMode_leading_square}.
\item \path{skewConstantMode_leading_coefficients}.
\end{itemize}

\subsection*{Group C5}\phantomsection\label{cert:C5}
\noindent Module: \path{XiFiveHalfCollisionOrbit.lean}.\\
Namespace: \path{RHTower}.
\begin{itemize}
\item \path{exists_xi_zero_five_half_three_half_ne_zero_of_not_mem_band}.
\item \path{xiFiveHalf_trapped_orbit_dichotomy}.
\end{itemize}

\subsection*{Group C6}\phantomsection\label{cert:C6}
\noindent Module: \path{ThetaCuspHawkingBMAngularSourceCollisionControl.lean}.\\
Namespace: \path{RHTower.D4BMJointSource}.
\begin{itemize}
\item \path{physicalSelfSourceCoefficient_ne_zero_of_four_fifths_le}.
\end{itemize}

\subsection*{Group C7}\phantomsection\label{cert:C7}
\noindent Module: \path{ThetaCuspHawkingTiltedBackreactionGaugeEndpoint.lean}.\\
Namespace: \path{RHTower}.
\begin{itemize}
\item \path{riemannHypothesis_iff_thetaCuspHawkingTiltedResonanceCentralizerRigid}.
\end{itemize}

\subsection*{Group C8}\phantomsection\label{cert:C8}
\noindent Module: \path{ThetaCuspHawkingKMSCayleyScatteringMixingEndpoint.lean}.\\
Namespace: \path{RHTower}.
\begin{itemize}
\item \path{riemannHypothesis_iff_thetaCuspHawkingKMSCayleyScatteringMixing}.
\end{itemize}

\subsection*{Group D2}\phantomsection\label{cert:D2}
\noindent Module: \path{ThetaCuspHawkingBMLiteralMixedActionCoefficient.lean}.\\
Namespace: \path{RHTower.D4BMMixedActionSource}.
\begin{itemize}
\item \path{literal_mixedPairedCoefficient_eq}.
\item \path{literal_affine_mixedPairedCoefficient}.
\end{itemize}

\subsection*{Group D3}\phantomsection\label{cert:D3}
\noindent Module: \path{ThetaCuspHawkingBMMixedActionCubicCoefficient.lean}.\\
Namespace: \path{RHTower.D4BMMixedActionSource}.
\begin{itemize}
\item \path{mixedPairedCoefficient_of_quadratic_zero}.
\end{itemize}

\subsection*{Group D4}\phantomsection\label{cert:D4}
\noindent Module: \path{ThetaCuspHawkingBMSourceContactRigidity.lean}.\\
Namespace: \path{RHTower.D4BMSourceContactRigidity}.
\begin{itemize}
\item \path{sourceContact_eq}.
\item \path{sourceContact_differential_identity}.
\item \path{sourceContact_rigidity_factor_ne_zero}.
\end{itemize}

\subsection*{Group D1}\phantomsection\label{cert:D1}
\noindent Module: \path{ThetaCuspHawkingBMPhysicalFundamentalFrequencies.lean}.\\
Namespace: \path{RHTower.D4BMJointSource}.
\begin{itemize}
\item \path{physicalFundamentalFrequency_eq}.
\item \path{physicalFundamentalFrequency_sub_sourceRate_ne_zero}.
\end{itemize}

\section{Explicit radial recipe for the three-field metric}\label{app:metric}

This appendix specifies the coefficients in \eqref{eq:metric}, so that the existence statement does not stand for an unspecified solution. All quantities depend on \(s,r\), \(r>1\); put \(d=(2s+1)(2s+2)\), \(f=r^2-1\), and use \(D_s\) as above. Define
\begin{align*}
 T&=\frac{s^2((s+1)^2-4r^2)}{8r^4},\\
 U_0&=r+\frac{s(3s+2)}{2r}+\frac{s^2(s+1)^2}{8r^3},&
 U_D&=\frac{3s(s+1)^2(3s+4)}{2r}.
\end{align*}
The two source rows required in the reconstruction are
\begin{align*}
 q_{tr}={}&\frac1{r^5f}\bigl[-s^2r^2+3s^2r^4+\tfrac34s^3
 -\tfrac72s^3r^2-\tfrac12s^3r^4+\tfrac{15}8s^4
 -\tfrac32s^4r^2\\
 &\hspace{33mm}+\tfrac32s^5-s^5r^2+\tfrac38s^6\bigr],\\
 q_t={}&\frac1{r^4}\bigl[\tfrac12sr^4+\tfrac12s^2r^2+\tfrac3{16}s^3
                  +\tfrac14s^3r^2+\tfrac38s^4+\tfrac3{16}s^5\bigr].
\end{align*}
They are \code{ssSourceTR}(s,r,0) and \code{ssSourceT}(s,r,0), respectively; both definitions are independent of the third argument. Now set
\begin{align*}
 H_0&=U_0/(4r),\\
 H_D&=\frac{(f/2)D_sU_D-dU_D/(4r)}{2-d},\\
 H_2&=\frac f2D_sU_0-\frac{U_0}{2r}
       -\frac{rf}{2s}q_{tr}-\frac fs q_t-\frac{U_D}{4r}+H_D,\\
 F_{10}&=0,\qquad F_{12}=\frac{4sr}{f}H_2,\\
 F_{1D}&=\frac{s}{f}(4rH_D-U_D).
\end{align*}
The remaining six coefficients are
\begin{align*}
 F_{20}&=-\frac{r}{2sf}q_{tr}+\frac{F_{1D}-F_{12}}{4srf}
                      -\frac{2H_0}{f^2}+\frac{2r}{f}D_sH_0,\\
 F_{22}&=\frac{F_{12}}{2srf}-\frac{2H_2}{f^2}+\frac{2r}{f}D_sH_2,\\
 F_{2D}&=\frac{dF_{1D}}{4srf}-\frac{2H_D}{f^2}+\frac{2r}{f}D_sH_D,\\
 F_{00}&=f^2F_{20}-2fT,\\
 F_{02}&=f^2F_{22}-2f(s+1)(s+2)T,\\
 F_{0D}&=f^2F_{2D}.
\end{align*}
The source definitions are in \path{ThetaCuspHawkingBMThreeFieldReconstruction.lean}, imported by group B3. The physical-rate identities exclude the denominators \(s=0\) and \(d=2\); \(r>1\) excludes the exterior coordinate denominators. These coefficients are used with the constructed angular fields, not with arbitrary functions merely assumed to solve the same equations.


\bibliographystyle{plain}
\bibliography{references}
\end{document}
